\begin{resumo}

JUNIOR, Edson Ramos da Silva; MOTTA, Bernardo Cordeiro. \textbf{Plataforma Web Escalável e Segura para Venda de Ingressos Online: Arquitetura, Requisitos e Implementação}. 2025. Plano de Projeto — Curso de Análise e Desenvolvimento de Sistemas, Centro Universitário de Viçosa, Viçosa, 2025.

A industria global de eventos enfrenta desafios significativos na comercializacao digital de ingressos, incluindo problemas de escalabilidade durante picos de demanda, vulnerabilidades de seguranca e experiencias de usuario inadequadas. Este trabalho apresenta o projeto arquitetural de uma plataforma web robusta para venda de ingressos online, fundamentada em principios de engenharia de software e centrada nas necessidades dos stakeholders. A metodologia adotada incluiu analise detalhada dos stakeholders atraves de personas e mapa de empatia, seguida pela especificacao completa de requisitos funcionais e nao funcionais. A solucao proposta contempla uma arquitetura full-stack baseada em Next.js e MongoDB com infraestrutura em nuvem escalavel, sistema de autenticacao seguro, gestao completa de eventos, fluxo de compra otimizado com integracao a gateways de pagamento, geracao de ingressos digitais com QR Codes dinamicos e sistema de controle de acesso com funcionamento offline. O projeto estabelece conformidade com padroes de seguranca PCI DSS, protecao contra vulnerabilidades OWASP Top 10, aderencia a LGPD e acessibilidade segundo WCAG 2.1. A analise de riscos e restricoes garante viabilidade tecnica e operacional da solucao, visando elevar os padroes de qualidade e confiabilidade no setor de eventos digitais.

\textbf{Palavras-chave:} sistemas de ingressos; arquitetura full-stack; Next.js; MongoDB; seguranca da informacao; escalabilidade.

\end{resumo}
